\documentclass[10pt]{extarticle}
\usepackage[a4paper,margin=1.5cm]{geometry}
\usepackage[T1]{fontenc}
\usepackage{amsmath,amssymb,mathtools}
\usepackage{graphicx}
\usepackage[export]{adjustbox}
\usepackage{multicol}
\usepackage{enumitem}
\usepackage{booktabs,array,tabularx}
\setlength{\columnsep}{0.7cm}
\setlength{\parindent}{0pt}
\setlength{\parskip}{3pt}
\pagestyle{empty}
\makeatletter
\newcommand{\qfigdef}[4]{\@namedef{qf@file@#1}{#2}\@namedef{qf@w@#1}{#3}\@namedef{qf@wide@#1}{#4}}
\newcommand{\qf@check}[1]{\@ifundefined{qf@file@#1}{\PackageError{score2}{Unknown figure id #1}{}}{}}
\newcommand{\qf@img}[4][]{\includegraphics[width=#2\dimexpr\@nameuse{qf@w@#4}pt\relax,max width=#3,max height=0.45\textheight,#1]{\@nameuse{qf@file@#4}}}
\newif\ifqf@full
\newif\ifqf@boxed
\newsavebox{\qf@box}
\newcommand{\qf@open}{\global\qf@fulltrue\par\global\qf@boxedtrue
  \setbox\qf@box\vbox\bgroup\hsize\textwidth\linewidth\textwidth}
\newcommand{\qf@emit}[1]{\par\egroup\global\qf@boxedfalse
  \ifdim\dimexpr\ht\qf@box+\dp\qf@box\relax<#1\textheight\noindent\usebox{\qf@box}\par\else\unvbox\qf@box\fi}
\newcommand{\qfullsplit}{\ifqf@boxed\qf@emit{0.85}\fi}
\newcommand{\qf@close}{\ifqf@boxed\qf@emit{0.45}\fi\global\qf@fullfalse}
\newcommand{\qfullbegin}{\end{multicols}\qf@open}
\newcommand{\qfullend}{\qf@close\begin{multicols}{2}}
\newcommand{\qfullstart}{\qf@open}
\newcommand{\qfulldone}{\qf@close}
\newcommand{\qfig}[2][1]{\qf@check{#2}%
  \ifqf@full\qf@blk{#1}{#2}\else
  \ifnum\@nameuse{qf@wide@#2}=1 \expandafter\@firstoftwo\else\expandafter\@secondoftwo\fi
  {\end{multicols}\begin{center}\qf@img{#1}{\textwidth}{#2}\end{center}\begin{multicols}{2}}%
  {\qf@blk{#1}{#2}}\fi}
\newcommand{\qf@blk}[2]{\par\smallskip{\centering\qf@img{#1}{\linewidth}{#2}\par}\smallskip}
\newcommand{\qfigi}[2][1]{\qf@check{#2}\qf@img[valign=c]{#1}{\linewidth}{#2}}
\makeatother
\newcommand{\papersec}[1]{\par\bigskip{\centering\Large\bfseries PAPER-#1\par}\medskip}
\newcommand{\partsec}[2]{\par\medskip\noindent{\large\bfseries PART-#1: #2}\par\smallskip}
\newcommand{\sect}[1]{\par\smallskip\noindent\textbf{#1}\par}
\newcommand{\qn}[1]{\par\smallskip\noindent\textbf{#1.}~}
\newcommand{\paraq}[2]{\par\smallskip\noindent\textbf{Paragraph for Questions #1 and #2}\par}
\newenvironment{qoptions}{\begin{itemize}[leftmargin=1.8em,itemsep=1pt,topsep=2pt,label={}]}{\end{itemize}}
\newcommand{\opt}[2]{\item[(#1)] #2}
\newcolumntype{L}[1]{>{\raggedright\arraybackslash}p{#1}}
\newcommand{\fwstart}{\end{multicols}\par\noindent\begin{minipage}{\textwidth}\small\setlength{\tabcolsep}{4pt}\renewcommand{\arraystretch}{1.5}}
\newcommand{\fwend}{\end{minipage}\par\medskip\begin{multicols}{2}}
% figure registry (generated by score2_tools.py finalize)
\qfigdef{p1_phy_s1_q01_1}{SRG (Enthusiast) - Test 04 (29 Nov 2020)-img/p1_phy_s1_q01_1.png}{109.5}{0}
\qfigdef{p1_phy_s1_q02_1}{SRG (Enthusiast) - Test 04 (29 Nov 2020)-img/p1_phy_s1_q02_1.png}{104.1}{0}
\qfigdef{p1_phy_s2_q05_1}{SRG (Enthusiast) - Test 04 (29 Nov 2020)-img/p1_phy_s2_q05_1.png}{339.9}{1}
\qfigdef{p1_phy_s2_q07_1}{SRG (Enthusiast) - Test 04 (29 Nov 2020)-img/p1_phy_s2_q07_1.png}{143.4}{0}
\qfigdef{p1_phy_s2_q08_1}{SRG (Enthusiast) - Test 04 (29 Nov 2020)-img/p1_phy_s2_q08_1.png}{45.6}{0}
\qfigdef{p1_phy_s2_q10_1}{SRG (Enthusiast) - Test 04 (29 Nov 2020)-img/p1_phy_s2_q10_1.png}{93.8}{0}
\qfigdef{p1_phy_s3_para11_13_1}{SRG (Enthusiast) - Test 04 (29 Nov 2020)-img/p1_phy_s3_para11_13_1.png}{175.2}{0}
\qfigdef{p1_phy_s4_q15_1}{SRG (Enthusiast) - Test 04 (29 Nov 2020)-img/p1_phy_s4_q15_1.png}{152.9}{0}
\qfigdef{p1_phy_s4_q18_1}{SRG (Enthusiast) - Test 04 (29 Nov 2020)-img/p1_phy_s4_q18_1.png}{144.5}{0}
\qfigdef{p1_phy_s4_q20_1}{SRG (Enthusiast) - Test 04 (29 Nov 2020)-img/p1_phy_s4_q20_1.png}{188.1}{0}
\qfigdef{p1_chem_s2_q09_1}{SRG (Enthusiast) - Test 04 (29 Nov 2020)-img/p1_chem_s2_q09_1.png}{137.5}{0}
\qfigdef{p1_chem_s2_q10_A}{SRG (Enthusiast) - Test 04 (29 Nov 2020)-img/p1_chem_s2_q10_A.png}{62.8}{0}
\qfigdef{p1_chem_s2_q10_B}{SRG (Enthusiast) - Test 04 (29 Nov 2020)-img/p1_chem_s2_q10_B.png}{103.2}{0}
\qfigdef{p1_chem_s2_q10_D}{SRG (Enthusiast) - Test 04 (29 Nov 2020)-img/p1_chem_s2_q10_D.png}{109.7}{0}
\qfigdef{p1_chem_s3_para11_13_1}{SRG (Enthusiast) - Test 04 (29 Nov 2020)-img/p1_chem_s3_para11_13_1.png}{110.4}{0}
\qfigdef{p1_chem_s3_q11_1}{SRG (Enthusiast) - Test 04 (29 Nov 2020)-img/p1_chem_s3_q11_1.png}{74.8}{0}
\qfigdef{p1_chem_s3_q11_2}{SRG (Enthusiast) - Test 04 (29 Nov 2020)-img/p1_chem_s3_q11_2.png}{58.0}{0}
\qfigdef{p1_chem_s3_q11_3}{SRG (Enthusiast) - Test 04 (29 Nov 2020)-img/p1_chem_s3_q11_3.png}{67.6}{0}
\qfigdef{p1_chem_s3_q11_4}{SRG (Enthusiast) - Test 04 (29 Nov 2020)-img/p1_chem_s3_q11_4.png}{68.9}{0}
\qfigdef{p1_chem_s3_q12_A}{SRG (Enthusiast) - Test 04 (29 Nov 2020)-img/p1_chem_s3_q12_A.png}{68.3}{0}
\qfigdef{p1_chem_s3_q12_B}{SRG (Enthusiast) - Test 04 (29 Nov 2020)-img/p1_chem_s3_q12_B.png}{75.7}{0}
\qfigdef{p1_chem_s3_q12_C}{SRG (Enthusiast) - Test 04 (29 Nov 2020)-img/p1_chem_s3_q12_C.png}{61.3}{0}
\qfigdef{p1_chem_s3_q12_D}{SRG (Enthusiast) - Test 04 (29 Nov 2020)-img/p1_chem_s3_q12_D.png}{60.6}{0}
\qfigdef{p1_chem_s3_q13_A}{SRG (Enthusiast) - Test 04 (29 Nov 2020)-img/p1_chem_s3_q13_A.png}{327.9}{1}
\raggedcolumns\makeatletter\@beginparpenalty=10000\makeatother

\begin{document}
\renewcommand{\paraq}[2]{\par\smallskip\noindent\textbf{Paragraph for Questions #1 to #2}\par}
\begin{center}{\Large\bfseries SRG (Enthusiast) --- Test 04 (29 Nov 2020)}\end{center}
\begin{multicols}{2}
\partsec{1}{PHYSICS}
\sect{SECTION-I (i)}
\qn{1} Consider two coherent monochromatic (wavelength $\lambda$) sources $S_1$ and $S_2$ separated by distance $d$. The ratio of intensity of $S_1$ and that of $S_2$ (which is responsible for interference at point P, where detector is located) is 4. The distance of point P from $S_1$ is (if the resultant intensity at point P is equal to $\dfrac94$ times of intensity of $S_1$)

[Given $\angle S_2S_1P$ is $90^\circ$ ($d>0$, $n$ is a positive integer)]
\qfig{p1_phy_s1_q01_1}
\begin{qoptions}
\opt{A}{$\dfrac{d^2-n^2\lambda^2}{2n\lambda}$}
\opt{B}{$\dfrac{d^2+n^2\lambda^2}{2n\lambda}$}
\opt{C}{$\dfrac{n\lambda d}{\sqrt{d^2-n^2\lambda^2}}$}
\opt{D}{$\dfrac{2n\lambda d}{\sqrt{d^2-n^2\lambda^2}}$}
\end{qoptions}
\qn{2} The refractive index of space changes with $y$, whose function is given as
\[\mu(y)=\begin{cases}1&\text{if }y\le0\\ ky^2&\text{if }0\le y\le2y_0\\ \dfrac{\ell}{y^2}&\text{if }2y_0\le y\le4y_0\end{cases}\]
where $k$ and $\ell$ are positive constant. When a light ray is incident at point $A$ as shown in figure. When the ray reaches to point $B$, it becomes parallel to $xz$-plane. Choose the correct statement.
\qfig{p1_phy_s1_q02_1}
\begin{qoptions}
\opt{A}{$\alpha$ is independent of $k$}
\opt{B}{$\alpha$ is independent of $\ell$}
\opt{C}{$\alpha$ is independent of both $k$ and $\ell$}
\opt{D}{None of these}
\end{qoptions}
\qn{3} The space has electromagnetic field which varies with time whose variation is given as
{\small\[\vec B=\begin{cases}-B_0\hat k&\text{if }0\le t\le\dfrac{\pi m}{qB_0}\\[5pt] -B_0\hat j&\text{if }\dfrac{\pi m}{qB_0}\le t\le\dfrac{2\pi m}{qB_0}\\[5pt] 0&\text{if }\dfrac{2\pi m}{qB_0}\le t\le\infty\end{cases}\]
and
\[\vec E=\begin{cases}0&\text{if }0\le t\le\dfrac{2\pi m}{qB_0}\\ -E_0\hat k&\text{if }\dfrac{2\pi m}{qB_0}\le t\le\infty\end{cases}\]}
A charge particle having mass $m$ and positive charge $q$ is given velocity $v_0\hat i$ at origin at $t=0$ sec. The coordinate of point on $xy$ plane when it is again passing through $xy$-plane for the first time is :
\begin{qoptions}
\opt{A}{$\left(\dfrac{2mv_0}{q}\sqrt{\dfrac{v_0}{E_0B_0}},\dfrac{2mv_0}{qB_0},0\right)$}
\opt{B}{$\left(\dfrac{2mv_0}{q}\sqrt{\dfrac{v_0}{E_0B_0}},\dfrac{mv_0}{qB_0},0\right)$}
\opt{C}{$\left(\dfrac{mv_0}{q}\sqrt{\dfrac{v_0}{E_0B_0}},\dfrac{2mv_0}{qB_0},0\right)$}
\opt{D}{$\left(\dfrac{mv_0}{q}\sqrt{\dfrac{v_0}{E_0B_0}},\dfrac{mv_0}{qB_0},0\right)$}
\end{qoptions}
\qn{4} A thin string with linear density $\mu$ is joined to a thick string with linear density $2\mu$. A incident pulse is sent down the thin string toward the thick string and eventually creates reflected and transmitted pulses. Which of following is true ?
\begin{qoptions}
\opt{A}{The reflected and transmitted pulses are both inverted}
\opt{B}{Neither the reflected nor transmitted pulses are inverted}
\opt{C}{The reflected pulse is inverted, but the transmitted pulse is not inverted}
\opt{D}{The transmitted pulse is inverted, but the reflected pulse is not inverted}
\end{qoptions}
\qfullbegin
\sect{SECTION-I (ii)}
\qn{5} In the shown figure (1) and (2), capacitors are in the steady state. Charging batteries are removed and switches $S_1$ and $S_2$ are closed at time $t=0$. The plot of $\ell n\,I$ ($I$ is the current in the resistor) in against time $t$ in the resistors $R_1$ and $R_2$ are shown by the graphs 1 and 2 respectively in the figure (3). Choose the correct option(s).
\qfig{p1_phy_s2_q05_1}
\qfullsplit
\begin{qoptions}
\opt{A}{If $\varepsilon_1=2\varepsilon_2$, $R_1$ must be equal to $2R_2$}
\opt{B}{If $R_1=R_2$, $C_1$ must be equal to $C_2$}
\opt{C}{If $C_1=C_2$, $R_1$ must be less than $R_2$}
\opt{D}{$R_1C_1$ is equal to $R_2C_2$}
\end{qoptions}
\qfullend
\qn{6} X-rays from a tube with a target $A$ of atomic number $Z$ shows strong $K$ line for target $A$ and weak $K$ lines for impurities. The wavelength of $K_\alpha$ lines is $\lambda_z$ for target $A$ and $\lambda_1$ and $\lambda_2$ for two impurities. If $\dfrac{\lambda_z}{\lambda_1}=4$ and $\dfrac{\lambda_z}{\lambda_2}=\dfrac14$ and taking screening constant of $K_\alpha$ lines to be unity, select the correct statement(s).
\begin{qoptions}
\opt{A}{The atomic number of first impurity is $2Z-1$}
\opt{B}{The atomic number of first impurity is $2Z+1$}
\opt{C}{The atomic number of second impurity is $\dfrac{(Z+1)}{2}$}
\opt{D}{The atomic number of second impurity is $\dfrac Z2+1$}
\end{qoptions}
\qn{7} A parallel plates capacitor is made of square conducting plates of side $a$ and the separation between plates is $d$. The capacitor is connected with battery of emf $V$ volt as shown in the figure. There is a dielectric slab of dimension $a\times a\times d$ with dielectric constant $k$. At $t=0$, dielectric slab is given velocity $v_0$ towards capacitor as shown in the figure. (Neglect the effect of gravity and electrostatic force acting on the dielectric when dielectric is out side of capacitor. Also ignore any type frictional force acting on the dielectric during its motion). Let the $x$ be the length dielectric inside the capacitor at $t=t$ sec. $\left[\ell_0\gg a\right]$
\qfig{p1_phy_s2_q07_1}
\begin{qoptions}
\opt{A}{Motion of dielectric slab is periodic but not simple harmonic motion}
\opt{B}{Motion of dielectric slab is simple harmonic motion}
\opt{C}{At any time, the slope of graph of total energy verses $x$ is twice the slope of graph of potential energy verses $x$.}
\opt{D}{The value of maximum energy stored in the system is $\dfrac12mv_0^2+\dfrac{\epsilon_0a^2V^2}{2d}(2k-1)$}
\end{qoptions}
\qn{8} A bob of mass $2m$ hangs by a string attached to the block of mass $m$ of a spring block system. The whole arrangement is in a state of equilibrium. The bob of mass $2m$ is pulled down slowly by a distance $x_0$ and released.
\qfig{p1_phy_s2_q08_1}
\begin{qoptions}
\opt{A}{For $x_0=\dfrac{3mg}{k}$, maximum tension in string is $4mg$}
\opt{B}{For $x_0=\dfrac{3mg}{k}$, minimum tension in string is $mg$}
\opt{C}{Frequency of oscillation of system is $\dfrac1{2\pi}\sqrt{\dfrac{k}{3m}}$, for all non-zero values of $x_0$}
\opt{D}{The motion will remain simple harmonic for $x_0\le\dfrac{3mg}{k}$}
\end{qoptions}
\qn{9} The equation of a process of diatomic gas is $P^2=\alpha^2V$ where $\alpha$ is a positive constant. Then choose the INCORRECT option(s).
\begin{qoptions}
\opt{A}{Work done by gas for a temperature change $T$ is $\dfrac{2}{3\alpha}nRT$}
\opt{B}{The change in internal energy is $\dfrac52nRT$ for a temperature change $T$}
\opt{C}{The specific heat for the process is $\dfrac{19}{9}R$}
\opt{D}{The change in internal energy for a temperature change $T$ is $\dfrac{5}{2\alpha}nRT$}
\end{qoptions}
\qn{10} A uniform rod $AB$ of length 7 m is undergoing combined rotational and translational motion on smooth horizontal surface such that, at some instant of time, velocities of its end point $A$ and centre $C$ are both perpendicular to the rod and opposite in direction, having magnitude 11 m/s and 3 m/s respectively as shown in the figure. Now consider that no external force and torque is acting on the rod. Then select CORRECT options.
\qfig{p1_phy_s2_q10_1}
\begin{qoptions}
\opt{A}{Magnitude of acceleration of point $A$ is 56 m/s$^2$}
\opt{B}{Magnitude of acceleration of point $B$ is 56 m/s$^2$}
\opt{C}{At the instant shown in the figure acceleration of point $C$ is zero}
\opt{D}{Angular velocity of the rod is 4 rad/s}
\end{qoptions}
\sect{SECTION-I (iii)}
\paraq{11}{13}
A block of mass 10 kg is put gently on a belt-conveyor system of infinite length at $t=0$ sec, which is moving with constant speed 20 m/sec rightward at all time, irrespectively of any situation by means of a motor-system as shown in the figure. A constant force of magnitude 15 N is applied on the block continuously during its motion.
\qfig{p1_phy_s3_para11_13_1}
\qn{11} The nature of friction force acting on the block of mass 10 kg during its motion is
\begin{qoptions}
\opt{A}{Kinetic only}
\opt{B}{Static only}
\opt{C}{Some time kinetic some time static}
\opt{D}{Can't be predicted}
\end{qoptions}
\qn{12} Work done by the kinetic friction on the block of mass 10 kg is
\begin{qoptions}
\opt{A}{1250 Joule}
\opt{B}{2500 Joule}
\opt{C}{$-1250$ Joule}
\opt{D}{None of these}
\end{qoptions}
\qn{13} The magnitude of acceleration of the block of mass 10 kg at $t=6$ sec is
\begin{qoptions}
\opt{A}{4 m/s$^2$}
\opt{B}{3 m/s$^2$}
\opt{C}{2 m/s$^2$}
\opt{D}{None of these}
\end{qoptions}
\sect{SECTION-II}
\qn{14} The pitch of a screw gauge is 1 mm and there are 50 divisions on its cap. When nothing is put in between the studs, 44$^{th}$ division of the circular scale coincides with the reference line and zero of main scale is not visible. When a plate is placed between the studs, the main scale reads 5 divisions and the circular scale reads 25 divisions. The true thickness of plate is $a\times10^{-5}$ m, find the value of $a$.
\qn{15} Two rings, each of mass $M=100$ gm are constrained to move along a fixed horizontal rod. An ideal string is connected with rings and a small $m=200$ gm is connected to the mid point of string. At any moment $m$ is moving downward with velocity $\sqrt3$ m/s. The speed of ring of $M$ at the same moment is
\qfig{p1_phy_s4_q15_1}
\qn{16} Pressure variation due to a sonic wave propagating along positive $x$-direction is given by the equation $\Delta P(x)=\Delta P_0\sin\left[2\pi(\upsilon t-x)\right]$. One end of an open organ pipe is closed by a cap and held at rest with its axis parallel to $x$-axis and the open end at the origin. What should be the minimum value of frictional force between the cap and the pipe so that pipe will remain closed. Given $\Delta P_0=2\times10^3$ N/m$^2$ and cross sectional area of the tube is equal to $5\times10^{-4}$ m$^2$,
\qn{17} In a slow reaction, heat is being evolved at a rate about 10 mW in a liquid. If the heat were being generated by the decay of $^{32}$P, a radioactive isotope of phosphorus that has half-life of 14 days and emits only beta-particles with a mean energy of 700 KeV, estimate the number of $^{32}$P atoms in the liquid. Express your answer in form $A\times10^{15}$ of and fill $A$ in answer sheet. [Take : $\ell n\,2=0.675$]
\qn{18} Shown in figure, a conical container of half-apex angle $37^\circ$ filled with certain quantities of kerosene and water. The force exerted by the water on the kerosene is $A\times10^5$ N fill value of $A$ in answer sheet, upto two decimal places. (Take atmospheric pressure, $10^5$ Pa)
\qfig{p1_phy_s4_q18_1}
\qn{19} The half-life of substance $X$ is 45 years, and it decomposes to substance $Y$. A sample from a meteorite was taken which contained 2\% of $X$ and 14\% of $Y$ by quantity of substance. If substance $Y$ is not normally found on a meteorite, what is the approximate age of the meteorite?
\qn{20} Water flowing along an open channel drives an undershot water wheel of radius 2 m (figure). The water approaches the wheel with a speed of 5.0 m/s and leaves with a speed of 2.5 m/s; the amount of water passing by is 20 kg per second. At what rate does the water deliver angular momentum to the wheel (in N-m) ?
\qfig{p1_phy_s4_q20_1}
\partsec{2}{CHEMISTRY}
\sect{SECTION-I (i)}
\qn{1} Among the following which one is more soluble in $\mathrm{H_2O}$
\begin{qoptions}
\opt{A}{CsF}
\opt{B}{RbF}
\opt{C}{NaF}
\opt{D}{LiF}
\end{qoptions}
\qn{2} Wave number of first line in Lyman series of hydrogen spectrum is equal to wave number of which one of the following lines:
\begin{qoptions}
\opt{A}{First line of Balmer series of $\mathrm{He^+}$}
\opt{B}{Second line of Paschen series of $\mathrm{Li^{2+}}$}
\opt{C}{Third line of Paschen series of $\mathrm{Li^{2+}}$}
\opt{D}{First line of Lyman series of $\mathrm{Li^{2+}}$}
\end{qoptions}
\qn{3} $\mathrm{CO_{2(g)}+C_{(s)}\rightleftharpoons2CO_{(g)}}$\quad $K_p=3$ bar at 727$^\circ$C. If initial pressure of $\mathrm{CO_2}$ is 0.48 bar. The partial pressure of $\mathrm{CO_2}$ at equilibrium is
\begin{qoptions}
\opt{A}{0.20 bar}
\opt{B}{0.30 bar}
\opt{C}{0.15 bar}
\opt{D}{0.33 bar}
\end{qoptions}
\qn{4} Which of the following elements has least electron affinity
\begin{qoptions}
\opt{A}{Tellurium}
\opt{B}{Sulphur}
\opt{C}{Oxygen}
\opt{D}{Selenium}
\end{qoptions}
\sect{SECTION-I (ii)}
\qn{5} Among the following salt solutions which give sulphide precipitate when $\mathrm{H_2S}$ gas is passed after the addition of $\mathrm{NH_4OH}$
\begin{qoptions}
\opt{A}{$\mathrm{Mn^{2+}}$}
\opt{B}{$\mathrm{Pb^{2+}}$}
\opt{C}{$\mathrm{Ni^{2+}}$}
\opt{D}{$\mathrm{Cu^{2+}}$}
\end{qoptions}
\qn{6} Aspirin and oil of winter green can be distinguished by
\begin{qoptions}
\opt{A}{Neutral $\mathrm{FeCl_3}$}
\opt{B}{aq. $\mathrm{NaHCO_3}$}
\opt{C}{aq.NaOH}
\opt{D}{$\mathrm{Br_2/CCl_4}$}
\end{qoptions}
\qn{7} Which of the following reaction/s is/are correct
\begin{qoptions}
\opt{A}{Monochlorination of 2-Methyl butane gives 2-Chloro-3-methylbutane as major product}
\opt{B}{Hydrogenation of cis-2-hexene gives a racemic mixture}
\opt{C}{Addition of excess HCl to acetylene gives ethyldenedichloride}
\opt{D}{Addition of HBr to cis-2-butene gives racemic mixture}
\end{qoptions}
\qn{8} Extraction of copper from copper pyrites involves
\begin{qoptions}
\opt{A}{Concentration of the powdered ore by froth flotation}
\opt{B}{Roasting of the ore in a reverbaratory furnace}
\opt{C}{Roasting of the copper-matte with silica to remove iron in the form of iron silicate slag}
\opt{D}{Leaching of the powdered ore with KCN in presence of air}
\end{qoptions}
\qn{9} Choose the correct statements with respect to following equilibrium
\qfig{p1_chem_s2_q09_1}
\begin{qoptions}
\opt{A}{\% of form I increases in water}
\opt{B}{\% of form I decreases in water}
\opt{C}{\% of form II increase in ether}
\opt{D}{\% of form I increases with increase in temperature in ether}
\end{qoptions}
\qn{10} Which of the following is/are able to exhibit optical isomerism?
\begin{qoptions}
\opt{A}{\qfigi{p1_chem_s2_q10_A}}
\opt{B}{\qfigi{p1_chem_s2_q10_B}}
\opt{C}{$\mathrm{CH_2=C=CH_2}$}
\opt{D}{\qfigi{p1_chem_s2_q10_D}}
\end{qoptions}
\sect{SECTION-I (iii)}
\paraq{11}{13}
A hydrocarbon $\mathrm{P(C_{11}H_8)}$ on treatment with $\mathrm{HgSO_4}$aq. $\mathrm{H_2SO_4}$ gives $\mathrm{Q(C_{11}H_{12}O_2)}$. P on treatment with $\mathrm{B_2H_6}$, THF and then with $\mathrm{H_2O_2}$/ NaOH gives R which is an isomer of Q. Q gives yellow precipitate with $\mathrm{I_2}$/NaOH, where as R do not give haloform reaction. R forms silver mirror on treatment with ammonical silver nitrate solution. P on boiling with alkaline $\mathrm{KMnO_4}$ solution forms $\mathrm{C_6H_6O_4}$, which on heating with $\mathrm{P_2O_5}$ gives $\mathrm{S(C_8H_4O_3)}$. S on treatment with toulene in presence of anhydrous $\mathrm{AlCl_3}$ gives $\mathrm{T\ (C_{15}H_{12}O_3)}$ as the major product. T is reduced by Zn-Hg/HCl and product obtained gives the following compound on reaction with HF.
\qfig{p1_chem_s3_para11_13_1}
\qn{11} Which of the following products are formed when Q is treated with aqueous NaOH?
\begin{qoptions}
\opt{a}{\qfigi{p1_chem_s3_q11_1}}
\opt{b}{\qfigi{p1_chem_s3_q11_2}}
\opt{c}{\qfigi{p1_chem_s3_q11_3}}
\opt{d}{\qfigi{p1_chem_s3_q11_4}}
\end{qoptions}
\begin{qoptions}
\opt{A}{a and d}
\opt{B}{a, b \& c}
\opt{C}{b and c}
\opt{D}{a, b, c and d}
\end{qoptions}
\qn{12} S on heating with hydrazine hydrochloride gives
\begin{qoptions}
\opt{A}{\qfigi{p1_chem_s3_q12_A}}
\opt{B}{\qfigi{p1_chem_s3_q12_B}}
\opt{C}{\qfigi{p1_chem_s3_q12_C}}
\opt{D}{\qfigi{p1_chem_s3_q12_D}}
\end{qoptions}
\qfullbegin
\qn{13} The structure of R is
\qfig{p1_chem_s3_q13_A}
\qfullend
\sect{SECTION-II}
\qn{14} Aqueous solutions of sodium iodate and sodium iodide react in acidic medium to give iodine, the total number of NaI molecules involved in the balanced chemical reaction is \underline{\hspace{1cm}}
\qn{15} A solution of 15.168 g of $\mathrm{Se_x}$ in 200 gm of benzene showed a depression in freezing point of 0.6$^\circ$C. If the molal depression constant of benzene is 5 K Kg mol$^{-1}$ and atomic weight of selenium is 79 u, the value of x is \underline{\hspace{1cm}}
\qn{16} Find out the value of expression $|x+y|$ for complex ion $\left[\mathrm{Co(EDTA)}\right]$. Where, $x$ = Total number of six membered ring of $y$ = Total number of five membered rings having oxygen donor atoms only.
\qn{17} Solubility of $\mathrm{Fe(OH)_2}$ ($k_{sp}=10^{-24}$) in a buffer solution containing 0.1 M BOH and 0.1 M $\mathrm{B^+}$ is $10^{-12}$. Calculate pKa of $\mathrm{B^{\oplus}}$.
\qn{18} All allotropic forms of boron contain $\mathrm{B_{12}}$ icosohedral units. If there are `p' faces, `q' corners and `r' edges in the icosohedron then $\dfrac{(p+q-r)}{2}$ is
\qn{19} The standard reduction potential of silver chloride electrode is 0.2 V and that of a silver electrode is 0.79 V. If the maximum amount of AgCl that can dissolve in $10^6$ Litres of a 0.1 M $\mathrm{AgNO_3}$ solution is $y$ millimoles, then the value of $y$ is \underline{\hspace{1cm}}
\qn{20} 0.2 moles of TlBr($k_{sp}=10^{-6}$) is added to 1 litre of a saturated solution of $\mathrm{Hg_2Br_2}$ ($k_{sp}=4\times10^{-24}$). The resultant solution has $2^x\times10^{-3y}$ M mercurous ion. The value of $(x+y)$ is
\partsec{3}{MATHEMATICS}
\sect{SECTION-I (i)}
\qn{1} Let $A=\begin{pmatrix}1&3&5\\4&2&6\\7&1&3\end{pmatrix}$; $B=\begin{pmatrix}3&6&7\\6&3&8\\7&8&3\end{pmatrix}$

Let $ABA^T=\begin{pmatrix}a_1&a_2&a_3\\a_4&a_5&a_6\\a_7&a_8&a_9\end{pmatrix}$

Then $(a_2-a_4)^2+(a_3-a_7)^2+(a_6-a_8)^2$ will be equal to
\begin{qoptions}
\opt{A}{0}
\opt{B}{13}
\opt{C}{78}
\opt{D}{65}
\end{qoptions}
\qn{2} $\displaystyle\int_1^5(x-1)(x-2)(x-3)(x-4)(x-5)\,dx$ lies in
\begin{qoptions}
\opt{A}{$[0,5]$}
\opt{B}{$[4,9]$}
\opt{C}{$[10,15]$}
\opt{D}{$[-5,0)$}
\end{qoptions}
\qn{3} $\displaystyle\int\dfrac{\left(3x^2+3x+1\right)dx}{x^6+3x^5+3x^4+x^3}=f(x)+k$. Then $f(1)=$ ?
\begin{qoptions}
\opt{A}{$-\dfrac18$}
\opt{B}{$-\dfrac78$}
\opt{C}{$-\dfrac38$}
\opt{D}{$-\dfrac58$}
\end{qoptions}
\qn{4} The minimum value of $\displaystyle\sum_{k=1}^5k\,|x-k|$ for all $x\in R$ is
\begin{qoptions}
\opt{A}{5}
\opt{B}{10}
\opt{C}{15}
\opt{D}{20}
\end{qoptions}
\sect{SECTION-I (ii)}
\qn{5} $\displaystyle\int_0^\pi\left(|\sin x|+|\cos x|+\big||\sin x|-|\cos x|\big|\right)dx$ equals to
\begin{qoptions}
\opt{A}{$\sqrt2$}
\opt{B}{$2\sqrt2$}
\opt{C}{$3\sqrt2$}
\opt{D}{$4\sqrt2$}
\end{qoptions}
\qn{6} Let $f(x)=\Big|\big||x|-3\big|-2\Big|-1\Big|$ where $x\in R$.

Let $a$ be the number of points where $f(x)$ is discontinuous

\hspace*{1.2em}$b$ be the number of points where $f(x)$ is non differentiable

\hspace*{1.2em}$c$ be the number of local minimas of $f(x)$

\hspace*{1.2em}$d$ be the number of local maximas of $f(x)$

Then $a+b+c+d$ is equal to
\begin{qoptions}
\opt{A}{22}
\opt{B}{24}
\opt{C}{26}
\opt{D}{28}
\end{qoptions}
\qn{7} The equation $\ell\mathrm n\,x=kx^4$ has atleast one solution. Then which of the following values of $k$ are possible
\begin{qoptions}
\opt{A}{$-\dfrac1e$}
\opt{B}{$\dfrac{1}{(4e)}$}
\opt{C}{$\dfrac{1}{(5e)}$}
\opt{D}{$\dfrac{1}{(10e)}$}
\end{qoptions}
\qn{8} The number of solutions of the equation $3x=100\{x\}$ is
\begin{qoptions}
\opt{A}{34}
\opt{B}{31}
\opt{C}{32}
\opt{D}{33}
\end{qoptions}
\qn{9} Consider a function $f:[-3,1]\to[-11,16]$ such that $f(x)=x^3-12x$. The function $f$ can be classified as
\begin{qoptions}
\opt{A}{One-one}
\opt{B}{Onto}
\opt{C}{Into}
\opt{D}{Many-one}
\end{qoptions}
\qn{10} Compute {\small\[\int_0^{10}\ell\mathrm n\left(x^7+1\right)dx+\int_0^{\ell\mathrm n(10000001)}\sqrt[7]{e^x-1}\,dx.\]}
\begin{qoptions}
\opt{A}{$\ell\mathrm n\left(10^7+1\right)$}
\opt{B}{$10\,\ell\mathrm n\left(10^7+1\right)$}
\opt{C}{$10^7\,\ell\mathrm n\left(10^7+1\right)$}
\opt{D}{$100\,\ell\mathrm n\left(10^7+1\right)$}
\end{qoptions}
\sect{SECTION-I (iii)}
\paraq{11}{13}
Sixteen players $S_1,S_2,\ldots.,S_{16}$ play in a tournament. They are divided into 8 pairs at random. From each pair a winner is decided on the basis of a game played between the two players of the pair. Assume that all the players are of equal strength.
\qn{11} The probability that exactly one of the two players $S_1$ \& $S_2$ is among the eight winners is
\begin{qoptions}
\opt{A}{$\dfrac4{15}$}
\opt{B}{$\dfrac7{15}$}
\opt{C}{$\dfrac8{15}$}
\opt{D}{$\dfrac9{15}$}
\end{qoptions}
\qn{12} What is the probability that all of $S_1,S_2,S_3,S_4$ are among the winners ?
\begin{qoptions}
\opt{A}{$\dfrac{1}{{}^{16}C_8}$}
\opt{B}{$\dfrac{{}^{12}C_4}{{}^{16}C_8}$}
\opt{C}{$\dfrac{{}^{8}C_4}{{}^{16}C_8}$}
\opt{D}{$\dfrac{4}{{}^{16}C_8}$}
\end{qoptions}
\qn{13} If the knockout procedure is repeated until a final winner emerges, what is the probability that $S_2$ wins the tournament ?
\begin{qoptions}
\opt{A}{$\dfrac1{16}$}
\opt{B}{$\dfrac18$}
\opt{C}{$\dfrac3{16}$}
\opt{D}{$\dfrac14$}
\end{qoptions}
\sect{SECTION-II}
\qn{14} If sum and the difference of the roots of the characteristic equation of a square matrix of order $2\times2$ are 4 and 2 respectively then the determinant of the matrix will be
\qn{15} Define functions $f_0,f_1,\ldots..,f_n,\ldots.$ to satisfy the following relations : $f_0=2$, $\dfrac{d}{dx}f_n=f_{n-1}$. If for $n\ge1$, $f_n(0)=0$, compute $\displaystyle\sum_{n=0}^{\infty}f_n(\ell\mathrm n(2019))$.
\qn{16} Let $A$ be the area of the region bounded by the curve $5x^2+6xy+2y^2+7x+6y+6=0$, then $[A]$ is equal to
\qn{17} If {\small$\begin{bmatrix}4a^2&4a&1\\4b^2&4b&1\\4c^2&4c&1\end{bmatrix}\begin{bmatrix}f(-1)\\f(1)\\f(2)\end{bmatrix}=\begin{bmatrix}3a^2+3a\\3b^2+3b\\3c^2+3c\end{bmatrix}$}, $f(x)$ is a quadratic function and its maximum value occurs at a point $V$. $A$ is a point of intersection of $y=f(x)$ with $x$-axis and point $B$ is such that chord $AB$ subtends a right angle at $V$. Let $A$ be the area enclosed by $f(x)$ and chord $AB$. Find $[A]$ (where $[\cdot]$ denotes the greatest integer function)
\qn{18} A tangent is drawn to the curve $y=x^3+x^2+x+1$ at $A$ where the curve intersects the line $y=4$. This tangent intersect the given curve at $B(x_1,\,y_1)$. Then $|\,x_1+y_1\,|$ is
\qn{19} Let $f:R\to R$ be a continuous and differentiable function. $f(x)=0$ has 10 real roots. Then the minimum number of real roots of the equation $f'\cdot f''+f\,f'''=0$ is
\qn{20} Consider the differential equation $y''-5y'+6y=0$. It can be converted to a linear differential equation of the form $y'-3y=c\cdot f(x)$. Then $f(\ell\mathrm n\,3)$ is
\end{multicols}
\end{document}
